\subsection{Bibliotecas de compresión sin pérdida consideradas}

Las alternativas evaluadas corresponden a bibliotecas utilizables desde C/C++ y
relevantes para compresión sin pérdida de datos residentes en memoria. Aunque sus
objetivos de diseño no son idénticos, todas representan opciones plausibles para
integrarse en una arquitectura de almacenamiento comprimido aplicada al vector de
estado.

\begin{itemize}
    \item \textbf{Blosc2:} biblioteca orientada a compresión de datos binarios y arreglos
    numéricos, organizada alrededor de contenedores particionados y de un pipeline de
    filtros y códecs configurable. Su diseño incorpora filtros como \emph{shuffle} y
    \emph{bitshuffle}, y permite trabajar con datos en memoria dentro de una
    organización por fragmentos \parencite{blosc2_docs}.

    \item \textbf{Zstandard (Zstd):} biblioteca generalista de compresión sin pérdida con
    un amplio rango de compromiso entre velocidad y ratio. Proporciona funciones de
    compresión y descompresión en memoria y resulta adecuada como códec base cuando
    el particionamiento y la gestión de bloques son responsabilidad de la aplicación
    \parencite{ZstdRepo, ZstdSite}.

    \item \textbf{LZ4:} biblioteca de compresión sin pérdida orientada a muy baja latencia,
    especialmente fuerte en descompresión. Su perfil resulta atractivo cuando el
    costo temporal del acceso tiene mayor peso que la tasa de compresión
    \parencite{LZ4Repo}.

    \item \textbf{zlib/Deflate:} solución madura y ampliamente adoptada para compresión
    sin pérdida en memoria y en flujo. Aunque no está especializada en arreglos
    científicos ni en acceso por bloques, constituye una referencia estable para
    contrastar bibliotecas más recientes \parencite{ZlibTechDetails}.
\end{itemize}

Las cuatro alternativas cubren perfiles distintos. Blosc2 representa una biblioteca con
orientación explícita a estructuras particionadas y datos numéricos en memoria; Zstd y
LZ4 son bibliotecas generalistas de alto desempeño que pueden emplearse como códecs
dentro de una arquitectura definida externamente; y zlib aporta una referencia clásica,
madura y ampliamente conocida.
